\section{Additional Model and Algorithm Details}
\label{supp:model-details}

\subsection{Detailed \(p\)-Fold Tucker Construction}
\label{supp:pfold-tucker}

Section~\ref{paper-sec:model} introduces the Tucker polynomial through
$p$-fold terms. Here we give the explicit index construction and show how sharing
the mode factors produces NTD-PL.

For a Tucker signal, a formal \(p\)-fold Tucker term may be written as
\begin{equation}
  \widetilde{Y}^{(p)}_{\mathbf{i}}
  =
  \sum_{\mathbf{r}^{(1)},\ldots,\mathbf{r}^{(p)}}
  \left(
  \prod_{\ell=1}^{p}
  X_{r_1^{(\ell)},\ldots,r_N^{(\ell)}}
  \right)
  \prod_{n=1}^{N}
  \mathcal{A}^{(n,p)}_{i_n,r_n^{(1)},\ldots,r_n^{(p)}} .
  \label{eq:supp-explicit-pfold-tucker}
\end{equation}
For \(p=1\), this is the ordinary Tucker entry formula. For \(p=2\), each
mode projection couples two Tucker index tuples, giving a quadratic
Tucker term. More generally, \eqref{eq:supp-explicit-pfold-tucker}
is the tensor analogue of augmenting a linear map with polynomial interactions
among Tucker coordinates.

For fixed mode maps, each output entry is a degree-$p$ polynomial in the core
entries. When $N=1$, this is an ordinary vector-valued homogeneous polynomial.
For $N>1$, the coefficient map is separated across the output modes in Tucker
form.

The general construction is not the parameterization used by the estimator:
each mode projection \(\mathcal{A}^{(n,p)}\) has \(I_nR_n^p\) entries and a
separate set of such projections would be needed for every active degree.
NTD-PL keeps the same interaction pattern but shares the original Tucker
factors across orders,
\begin{equation}
  \mathcal{A}^{(n,p)}_{i_n,r_n^{(1)},\ldots,r_n^{(p)}}
  =
  \prod_{\ell=1}^{p} A^{(n)}_{i_n,r_n^{(\ell)}} .
  \label{eq:supp-shared-projection-tying}
\end{equation}
Substituting \eqref{eq:supp-shared-projection-tying} into
\eqref{eq:supp-explicit-pfold-tucker} gives exactly
\(\cS^{\od p}\), where
\(\cS=\tucker{\cX}{\bA^{(1)},\ldots,\bA^{(N)}}\). Thus a degree-\(p\)
NTD-PL term is a shared-factor $p$-fold Tucker term: it couples $p$ Tucker
index tuples while reusing the same mode factors.

This viewpoint explains why the model can increase effective rank without
storing an untied high-order tensor block. The
additional trainable variables are the response coefficients, while the
high-order features are induced by repeated use of the same Tucker factors.

\subsection{Projection Interpretation of the Response Diagnostic}

The diagnostic $\Dlink$ in \eqref{paper-eq:dlink} can be interpreted as a projection
test for residual structure that depends on the current Tucker estimate. Let
$\bPhi_{\Omega,J}$ be the response feature matrix built from a fixed Tucker
reconstruction $\cS_T$ and let $\mathbf{P}_J$ denote the orthogonal projector
onto its column space. When $\lambda_\beta=0$, the refreshed-response residual
satisfies
\begin{equation}
  R_{\mathrm{resp}}
  =
  \norm{(\bI-\mathbf{P}_J)y_\Omega}_2^2.
\end{equation}
Thus, $\Dlink$ is the residual-energy ratio, in dB, between the identity
response and the projection onto the response-feature subspace generated by the
fixed Tucker fit. With ridge regularization, the response coefficients are
obtained from the corresponding shrinkage operator, while $R_{\mathrm{resp}}$
still denotes the data residual without the ridge penalty. A large positive
$\Dlink$ means that the residual is aligned with scalar functions of the current
Tucker estimate, whereas a small or nonpositive $\Dlink$ suggests increasing
multilinear rank, changing the tensor prior, or revisiting the observation model.

\subsection{Masked-Entry Generalization}

Let \(\Theta\subset\R^d\) be a bounded parameter set for a degree-\(P\)
NTD-PL model, where
\begin{equation}
  d=
  \prod_{n=1}^{N}R_n
  +
  \sum_{n=1}^{N}I_nR_n
  +
  (P+1).
\end{equation}
For \(\theta\in\Theta\), define the all-entry risk
\(\mathcal{L}(\theta)=M^{-1}\sum_{\mathbf{i}}\ell_{\mathbf{i}}(\theta)\),
\(M=\prod_n I_n\), and the masked empirical risk
\(\widehat{\mathcal{L}}_m(\theta)=m^{-1}\sum_{t=1}^{m}
\ell_{\mathbf{i}_t}(\theta)\), where the entries
\(\mathbf{i}_t\) are sampled uniformly.

\begin{theorem}[Masked-entry uniform convergence]
\label{thm:supp-masked-generalization}
Assume \(0\le \ell_{\mathbf{i}}(\theta)\le B_\ell\) and
\(\ell_{\mathbf{i}}(\theta)\) is \(G\)-Lipschitz in \(\theta\), uniformly over
\(\mathbf{i}\). If \(\Theta\) lies in a Euclidean ball of radius \(B_\Theta\),
then, with probability at least \(1-\delta\),
\begin{equation}
  \sup_{\theta\in\Theta}
  \left|
  \mathcal{L}(\theta)-\widehat{\mathcal{L}}_m(\theta)
  \right|
  \le
  2G\varepsilon
  +
  B_\ell
  \sqrt{
  \frac{
  d\log(3B_\Theta/\varepsilon)+\log(2/\delta)}
  {2m}}
\end{equation}
for every \(0<\varepsilon\le B_\Theta\).
\end{theorem}

This standard covering bound identifies the complexity entering masked
recovery: the dimension is the shared-factor parameter count, not
that of an untied Tucker model with Veronese ranks.

\subsection{Detailed Computational Scaling}

Let $M=\prod_{n=1}^{N}I_n$, $R_\Pi=\prod_{n=1}^{N}R_n$, and
$d_{\mathrm{Tucker}}=R_\Pi+\sum_{n=1}^{N}I_nR_n$. Storing the
model requires
\begin{equation}
  d_{\mathrm{NTD}}=d_{\mathrm{Tucker}}+(J+1)
\end{equation}
parameters. For the power basis, $J=P$. A higher-rank Tucker model that
explicitly matches the degree-$P$ Veronese lift may require mode ranks
$D_n=\binom{R_n+P-1}{P}$, with core size $\prod_nD_n$ and factor size
$\sum_nI_nD_n$. NTD-PL avoids storing this larger Tucker tensor.

Response Refresh forms $\bPhi_{\Omega,J}^{T}\bPhi_{\Omega,J}$ and solves a
small system, costing $\mathcal{O}(mJ^2+J^3)$. Dense Tucker reconstruction and
reverse contractions scale as
\begin{equation}
  \mathcal{O}\!\left(
  M R_\Pi + \sum_{n=1}^{N} I_n R_n R_\Pi/R_n
  \right)
\end{equation}
up to contraction-order constants, followed by $\mathcal{O}(MJ)$ work for the
entrywise response. In masked settings, response and residual computations use
$m$ observed entries. Since $J$ is small, the response update is typically
minor relative to the Tucker contractions.
